\subsection{Introduction \quad / \quad \textthai{บทนำ}}
\begin{paracol}{2}
Modern physics research is increasingly becoming a data-intensive endeavor. Whether analyzing petabytes of data from particle colliders or simulating the evolution of galaxies, the choice of tools determines the efficiency and reproducibility of the science. Python has emerged as the \textit{lingua franca} of scientific computing due to its readability, vast ecosystem, and high-performance libraries that wrap C and Fortran code.

In this chapter, we delve into the core libraries that form the "Physicist's Python Toolkit": \texttt{NumPy}\index{NumPy@NumPy (NumPy)}, \texttt{SciPy}\index{SciPy@SciPy (SciPy)}, \texttt{Matplotlib}, and \texttt{Pandas}. Mastering these is a prerequisite for the more advanced AI techniques discussed in later chapters.

\switchcolumn

\begin{thai}
การวิจัยทางฟิสิกส์สมัยใหม่ได้กลายเป็นความพยายามที่ต้องใช้ข้อมูลอย่างมหาศาลมากขึ้นเรื่อยๆ ไม่ว่าจะเป็นการวิเคราะห์ข้อมูลระดับเพตาไบต์จากเครื่องเร่งอนุภาค หรือการจำลองวิวัฒนาการของกาแล็กซี การเลือกใช้เครื่องมือเป็นตัวกำหนดประสิทธิภาพและความสามารถในการทำซ้ำของวิทยาศาสตร์ ภาษา Python ได้กลายเป็น "ภาษากลาง (lingua franca)" ของการคำนวณทางวิทยาศาสตร์ เนื่องจากความง่ายในการอ่าน ระบบนิเวศที่กว้างขวาง และไลบรารีประสิทธิภาพสูงที่พัฒนาต่อยอดมาจากรหัส C และ Fortran

ในบทนี้เราจะเจาะลึกไลบรารีหลักที่เป็น "ชุดเครื่องมือ Python สำหรับนักฟิสิกส์": \texttt{NumPy}, \texttt{SciPy}, \texttt{Matplotlib} และ \texttt{Pandas} การเชี่ยวชาญเครื่องมือเหล่านี้เป็นข้อกำหนดเบื้องต้นสำหรับเทคนิค AI ขั้นสูงที่จะกล่าวถึงในบทต่อๆ ไป
\end{thai}
\end{paracol}
